\section{Results}

\subsection{Results After Applying the Approach}

Our research question is: \textbf{How can we accurately identify the implementation of a specific business rule from its textual description?} We address this question incrementally. At its current stage, our approach does not aim to directly pinpoint the exact code location implementing a business rule. Rather, it is designed to support developers by narrowing the search space to a manageable set of candidate methods. More specifically, the approach filters out methods that are unlikely to be related to the business rule while retaining those that may potentially implement it. Our objective is therefore to assess the extent to which the approach can reduce the search space and demonstrate its practical relevance for business rule localization.

We use two complementary metrics to assess the effectiveness of our approach. The first metric, \textbf{Success Rate}, corresponds to \textbf{Recall} and measures the proportion of comments for which the method associated with the business rule is successfully retained by the approach. In other words, it evaluates the ability of the approach to retain the method implementing the business rule among the candidate methods. The second metric, \textbf{Search Space Reduction}, measures the proportion of methods eliminated by the approach, thereby quantifying the extent to which the original search space is reduced.

We do not include \textbf{Precision} among our evaluation metrics, as its strict interpretation is not well suited to our objective. Precision strongly penalizes cases in which multiple candidate methods are returned and can decrease substantially even when only a small number of additional methods are considered. For example, returning two candidate methods instead of a single one would reduce the precision by half. However, from a developer's perspective, examining two candidate methods is still a simple and time-efficient task. Therefore, a low precision value does not necessarily reflect poor practical performance in our context.

Our approach is not intended to directly identify a single method, but rather to reduce the search space to a manageable set of candidate methods. Consequently, we focus on whether the resulting set is sufficiently small to be efficiently investigated by a developer. The objective is thus not to minimize the number of returned methods to exactly one, but to provide a sufficiently restricted and manageable set of candidates that facilitates the developer's investigation.

\subsubsection{Success Rate}\noindent

To evaluate our approach, we applied it to a dataset comprising 100 textual descriptions of business rules. For each description, we examined whether the method associated with the corresponding business rule was retained among the candidate methods returned by the approach. The associated method was successfully retained in 96 cases (96\%), demonstrating that the approach was able to preserve the relevant method for the vast majority of the evaluated business rules. In the remaining 4 cases (4\%), the associated method was not included in the final set of candidate methods. These results indicate that the approach achieves a high recall while substantially narrowing the search space.

\subsubsection{Search Space Reduction}\noindent

We further analyzed the 96 cases in which the method associated with the business rule was successfully retained by our approach. This analysis focuses on the size of the resulting candidate sets and, more importantly, on the extent to which the approach reduces the developer's search space.

Across these 96 cases, the number of candidate methods returned by the approach ranged from 2 to 1,248, with a median of 139 methods. Although some business rule descriptions resulted in relatively large candidate sets, half of the cases resulted in fewer than 139 candidate methods. These results provide an indication of the extent to which the approach narrows the search space available to developers.

However, the absolute number of retained methods should not be considered in isolation. Our objective is not to directly identify a single method implementing a business rule, but rather to substantially reduce the number of methods that a developer needs to investigate. We therefore consider the proportion of methods eliminated by the approach to be a more meaningful indicator of its practical contribution.

In the studied project, the initial codebase contains 115,829 methods. After filtering this set and retaining only the methods that are more likely to contain business logic, we reduced the number of candidate methods to 14,578. We then applied our approach to these methods to generate a candidate set for each business rule description.

For the 96 successfully retained cases, the approach eliminated the vast majority of the 14,578 candidate methods. Even in the worst-case scenario, where 1,248 methods were retained, the approach eliminated more than 91\% of the candidate methods. This corresponds to a reduction from 14,578 methods to only 1,248 candidates. In most cases, the reduction was even more substantial.

These results demonstrate that, despite not necessarily identifying a single implementation method, our approach can considerably narrow the search space. By reducing the number of methods that need to be manually investigated, it provides developers with a more manageable set of candidates and facilitates the subsequent identification of the method implementing the business rule.

\subsection{Results of the Context-Based Refinement}

Applying the Context-Based Refinement substantially reduces the size of the candidate method sets while maintaining a high recall. Before applying the refinement, the number of candidate methods ranged from 2 to 1,248, with a median of 139 methods per business rule description. After applying the context-based filter, these values decreased to a maximum of 179 methods, a minimum of 2, a median of 18.5 methods. This represents a substantial reduction in the size of the candidate sets. In particular, the worst-case search space reduction increased from 91\% to 99\%, meaning that even in the case with the largest number of retained methods, 99\% of the initial candidate methods were eliminated. This reduction is achieved while maintaining a high recall, which decreases only slightly from 96\% to 94\%. These results indicate that the contextual information provided by the document and section names can effectively refine the candidate set and considerably reduce the effort required for the developer to investigate the remaining methods.